\documentclass[10pt,a4paper]{amsart}
\usepackage[margin=19mm]{geometry}
\usepackage{amsmath,amssymb,mathtools}
\usepackage[scr=rsfs,cal=euler]{mathalfa}
\usepackage{xfrac}
\usepackage[unicode,hidelinks,bookmarks=false]{hyperref}
\newcommand{\ie}{i.e.\ }
\newcommand{\cf}{cf.\ }
\newcommand{\resp}{respectively\ }
\newcommand{\Subsection}{\subsection}
\providecommand{\nobreakdash}{\nobreak\hbox{-}}
\setlength{\emergencystretch}{2em}
\multlinegap0pt
\usepackage{amssymb}
\usepackage[shortlabels]{enumitem}
\usepackage{mathtools}
\usepackage{tikz}
\usepackage[normalem]{ulem}
\usepackage{xcolor}
\newtheorem{lem}{Lemma}
\newtheorem{thm}{Theorem}
\title{On the Ehrhart Theory of Generalized Symmetric Edge Polytopes}
\author{Robert Davis, Akihiro Higashitani, Hidefumi Ohsugi}
\date{Source: arXiv:2401.03383v4 (CC BY 4.0)}
\begin{document}
\maketitle

\noindent\emph{Source and licence.} On the Ehrhart Theory of Generalized Symmetric Edge Polytopes. Version arXiv:2401.03383v4 (CC BY 4.0), distributed by arXiv under the Creative Commons Attribution 4.0 International licence, http://creativecommons.org/licenses/by/4.0/, as recorded in the arXiv licence metadata for this exact version and shown on the abstract page https://arxiv.org/abs/2401.03383v4. Full licence notice: \texttt{licenses/arxiv-2401.03383-CC-BY-4.0.txt}.\par\smallskip

\noindent\textit{Editorial scope.} Counted unit: the article's thm gamma (source0.tex lines 1809--1813). The article's complete original proof of it is delivered with it (source0.tex lines 1816--2197, one \texttt{proof} environment of 382 lines). No mathematics is written, repaired or re-derived: the statement and the proof are the article's own lines, in the article's own notation, and no retained line is substituted or re-flowed.  Retained uncounted as the source-local context the proof needs: lines 1662--1667; lines 1710--1715; lines 1767--1772.  Nothing was omitted from any retained span, so the package carries no omission record.  No retained line contains a citation, so the wrapper carries no bibliography and no bibliography entry.  No retained line contains a figure environment, an image-inclusion command or drawing code, so no image is delivered and the package declares no asset dependency.  Editorial formatting only, in addition: an AMS \texttt{amsart} preamble that restates the article's own declarations and macros used by the retained lines, namely the environments \texttt{lem}, \texttt{thm} on separate counters, together with the standard packages those lines need.  The retained environments print in this document as Lemma 1, Theorem 1 in order of appearance, whereas the article prints the same items with the numbering of its own full document.  The complete native source of the article as distributed in the arXiv e-print for this exact version is bundled unchanged as \texttt{sources/152574/source0.tex}.
\begin{lem}\label{hodai1}
    For all $0 \leq k \leq \ell$ we have
    \[
        \sum_{a = 0}^\ell(-1)^{k-a}4^{\ell-a} \binom{2a}{a} \binom{\ell-k}{\ell-a} = \frac{\binom{\ell}{k} \binom{2\ell}{\ell}}{ \binom{2\ell}{2k}}.
    \]    
\end{lem}

\begin{lem} \label{hodai2}
For all $k \geq 0$,
\[
    \binom{2k}{k}\frac{x^k}{(1-4x)^{k+\frac{3}{2}}} = \sum_{\ell \ge 0} \frac{2\ell+1}{2k+1}\binom{\ell}{k}\binom{2\ell}{\ell}x^\ell.
\]  
\end{lem}

\begin{lem} \label{hodai3}
For $\ell \geq 0$,
\[
    \sum_{p,q \ge 0}\binom{2 \ell +1}{p+q+1} x^p y^q = \frac{(x+1)^{2\ell+1} - (y+1)^{2\ell+1}}{x-y}.
\]
\end{lem}

\begin{thm}\label{thm: f to gamma}
For all $\ell \geq 0$ we have
    \[\sum_{p+q \le 2 \ell} \binom{2 \ell +1}{p+q+1}t^{2\ell -p-q} \sum_{i=0}^p \sum_{j=0}^q  \binom{p}{i} \binom{q}{j}\binom{2\ell -p-q}{\ell -q-i+j} t^{2i + 2j} = \sum_{a=0}^\ell \binom{2a}{a} t^a (t+1)^{4\ell -2a}.
    \]
\end{thm}

\begin{proof}
For notational convenience, set 
\[
F_\ell = \sum_{a=0}^\ell \binom{2a}{a} t^a (t+1)^{4\ell -2a}
\]
and
\[
G_{\ell,m} = 
\sum_{p+q \le 2 \ell} \binom{2 \ell +1}{p+q+1}
t^{2\ell -p-q} \sum_{i=0}^p \sum_{j=0}^q  \binom{p}{i} \binom{q}{j}
\binom{2m -p-q}{m -q-i+j} t^{2i + 2j}.\]
We compute the generating function for $F_{\ell}$ as follows:
\begin{eqnarray*}
\sum_{\ell \ge 0}  F_\ell  \lambda^\ell
&=&\sum_{\ell \ge 0} 
\sum_{a=0}^\ell \binom{2a}{a} t^a (t+1)^{4\ell -2a}
\lambda^\ell\\
&=&
\sum_{a\ge 0} 
\binom{2a}{a} t^a 
\sum_{\ell \ge a} 
(t+1)^{4\ell -2a} 
\lambda^\ell\\
&=&
\sum_{a\ge 0} 
\binom{2a}{a} t^a 
\frac{(t+1)^{2a} \lambda^a}{
1- (t+1)^4 \lambda
}\\
&=&
\frac{1}{
1- (t+1)^4 \lambda
}\sum_{a\ge 0} 
\binom{2a}{a} t^a 
(t+1)^{2a} \lambda^a\\
&=&
\frac{1}{
(1- (t+1)^4 \lambda)\sqrt{
1-4 t (t+1)^2 \lambda
}}.
\end{eqnarray*}
Hence it is enough to show that
\[
\sum_{\ell \ge 0} G_{\ell,\ell} \lambda^\ell
=\frac{1}{
(1- (t+1)^4 \lambda)\sqrt{
1-4 t (t+1)^2 \lambda
}}.
\]
Using Lemma \ref{hodai3}, we have
\[\begin{aligned}
\sum_{m \ge 0} G_{\ell, m} \lambda^m
&=
%\sum_{\ell \ge 0}
\sum_{m \ge 0}
\sum_{p, q \ge 0}
\binom{2 \ell +1}{p+q+1}
t^{2\ell -p-q} \sum_{i=0}^p \sum_{j=0}^q  \binom{p}{i} \binom{q}{j}
\binom{2 m -p-q}{m -q-i+j} t^{2i + 2j}
\lambda^m\\
&=
\sum_{p,q \ge 0}
\sum_{i=0}^p \sum_{j=0}^q  
\binom{p}{i} \binom{q}{j}
t^{2\ell -p-q + 2i + 2j}
%
%\sum_{\ell \ge 0}
\binom{2 \ell +1}{p+q+1}
%\lambda^\ell
%
\sum_{m \ge 0}
\binom{2m -p-q}{m -q-i+j} 
\lambda^m
\\
&=
\sum_{p,q \ge 0}
\sum_{i=0}^p \sum_{j=0}^q  
\binom{p}{i} \binom{q}{j}
t^{2\ell -p-q + 2i + 2j}
%
%\sum_{\ell \ge 0}
\binom{2 \ell +1}{p+q+1}
%\lambda^\ell
\lambda^{q + i -j}
%
\sum_{m \ge 0}
\binom{2m -p +q + 2i -2j}{m} 
\lambda^{m}
\\
&=
\sum_{p,q \ge 0}
\sum_{i=0}^p \sum_{j=0}^q  
\binom{p}{i} \binom{q}{j}
t^{2\ell -p-q + 2i + 2j}
%
%\sum_{\ell \ge 0}
\binom{2 \ell +1}{p+q+1}
%
\frac{ \lambda^{q + i -j}}{\sqrt{1 - 4 \lambda}}
%\left(
%\frac{1-\sqrt{1 - 4 \lambda}}{2 \lambda}
%\right)
\ \ L^{-p +q + 2i -2j}
\\
&=
\frac{t^{2\ell}}{\sqrt{1 - 4 \lambda}}
\sum_{p,q \ge 0}
\binom{2 \ell +1}{p+q+1}
\left(
\lambda L t + \frac{1}{tL}
\right)^p
\left(
\frac{\lambda L}{t}+\frac{t}{L}
\right)^q
\\
&=
\frac{t^{2\ell}}{\sqrt{1 - 4 \lambda}}
\frac{
(\lambda L t +\frac{1}{tL}+1)^{2\ell+1}
-
(\frac{\lambda L}{t}+\frac{t}{L}+1)^{2\ell+1}
}
{
(\lambda L t + \frac{1}{tL})-(\frac{\lambda L}{t}+\frac{t}{L})
},
\end{aligned}
\]
where $L =\frac{1-\sqrt{1 - 4 \lambda}}{2 \lambda}$.
%Note that $M$ satisfies $\lambda M^2-M+1=0$.
It follows that $\frac{1}{L} = \frac{1+\sqrt{1 - 4 \lambda}}{2}$, that
\[\begin{aligned}
    \lambda L t +\frac{1}{tL}+1 &=  \frac{1-\sqrt{1 - 4 \lambda}}{2} t
+\frac{1+\sqrt{1 - 4 \lambda}}{2t } +1\\
&=
\frac{1}{2t} 
\left(\sqrt{1 - 4 \lambda} (1-t^2) +(t+1)^2 \right)\\
&=
\frac{t+1}{2t} 
\left(\sqrt{1 - 4 \lambda} (1-t) +(t+1) \right),
\end{aligned}\]
and that
\[\begin{aligned}
\frac{\lambda L}{t}+\frac{t}{L}+1
&=\frac{1-\sqrt{1 - 4 \lambda}}{2t}
+\frac{1+\sqrt{1 - 4 \lambda}}{2} t +1\\
&=
\frac{1}{2t} 
\left(\sqrt{1 - 4 \lambda} (t^2-1) + (t+1)^2 \right)\\
&=
\frac{t+1}{2t} 
\left(\sqrt{1 - 4 \lambda} (t-1) +(t+1) \right).
\end{aligned}
\]
Hence,
\[
%
%
%
\left(\lambda L t +\frac{1}{tL}\right)-\left(\frac{\lambda L}{t}+\frac{t}{L}\right) = 
\left(\lambda L t +\frac{1}{tL}+1\right)-\left(\frac{\lambda L}{t}+\frac{t}{L}+1\right) =
\frac{1-t^2}{t}
\sqrt{1-4\lambda}.
\]
Thus
\[
\begin{aligned}
    \sum_{m \ge 0} G_{\ell, m} \lambda^m
%&=&
%\frac{1}{(1-t^2)(1 - 4 \lambda)}
%\\
%& & \times
%\left(
%\left(\frac{(t+1) \left(
%2 \lambda (1-t) +t (1-\sqrt{1-4 \lambda})
%\right)}{1-\sqrt{1-4 \lambda}}
%\right)^{2\ell+1}
%-
%\left(\frac{(t+1) \left(2 \lambda( t-1)
%+1-\sqrt{1-4 \lambda}\right)}{1-\sqrt{1-4 \lambda}}
%\right)^{2\ell+1}
%\right)\\
%&=&
%\frac{1}{(1-t^2) (1 - 4 \lambda)}
%\left(
%\left(\frac{2 \lambda (1-t^2)}{1-\sqrt{1 - 4 \lambda}} +t(t+1)  \right)^{2\ell+1}
%-
%\left(\frac{2 \lambda (t^2-1)}{1-\sqrt{1 - 4 \lambda}} +(t+1) \right)^{2\ell+1}
%\right)\\
%&=&
%\frac{1}{(1-t^2) (1 - 4 \lambda)}
%\left(
%\left(
%\frac{\sqrt{1 - 4 \lambda} (1-t^2)+(t+1)^2}{2}  
%\right)^{2\ell+1}
%+
%\left(
%\frac{\sqrt{1 - 4 \lambda} (1-t^2) -(t+1)^2}{2} 
%\right)^{2\ell+1}\right)\\
&=
\frac{(t+1)^{2\ell +1}}{2^{2 \ell +1} (1-t^2) (1 - 4 \lambda)}
\left(
\left(
\sqrt{1 - 4 \lambda} (1-t)+(t+1)
\right)^{2\ell+1}
+
\left(
\sqrt{1 - 4 \lambda} (1-t) -(t+1)
\right)^{2\ell+1}\right)\\
&=
\frac{(t+1)^{2\ell +1}}{2^{2 \ell} (1-t^2) (1 - 4 \lambda)}
\sum_{k\ge 0}
\binom{2\ell+1}{2k}
(\sqrt{1 - 4 \lambda} (1-t))^{2\ell+1-2k}
(t+1)^{2k}
\\
&=
\frac{(t+1)^{2\ell} (1-t)^{2\ell} }{2^{2 \ell} \sqrt{1 - 4 \lambda} }
\sum_{k\ge 0}
\binom{2\ell+1}{2k}
(1-4\lambda)^{\ell-k}
\left(
\frac{t+1}{1-t}
\right)
^{2k}
\\
&=
4^{-\ell} (t+1)^{2\ell} (1-t)^{2\ell}\sum_{a \ge 0}
\binom{2a}{a} \lambda^a
\sum_{k\ge 0}
\binom{2\ell+1}{2k}
(1-4\lambda)^{\ell-k}
\left(
\frac{t+1}{1-t}
\right)
^{2k}
\\
&=
4^{-\ell} (t+1)^{2\ell} (1-t)^{2\ell}
\sum_{a \ge 0}
\binom{2a}{a} \lambda^a
\sum_{k\ge 0}
\binom{2\ell+1}{2k}
\sum_{i=0}^{\ell-k}
\binom{\ell-k}{i}
(-4\lambda)^{i}
\left(
\frac{t+1}{1-t}
\right)
^{2k}.
\end{aligned}\]
Since $G_{\ell, \ell}$ is the coefficient of $\lambda^\ell$ in $\sum_{m \ge 0} G_{\ell, m} \lambda^m$,
we have
\[
G_{\ell, \ell} = 
4^{-\ell} (t+1)^{2\ell} (1-t)^{2\ell}\sum_{a \ge 0}
\sum_{k\ge 0}
\binom{2a}{a} 
\binom{2\ell+1}{2k}
\binom{\ell-k}{\ell-a}
(-4)^{\ell-a}
\left(
\frac{t+1}{1-t}
\right)
^{2k}.
%\\
%&=& 
%(-1)^\ell
%(t+1)^{2\ell} (1-t)^{2\ell}
%\sum_{a \ge 0}
%\sum_{k=0}^\ell
%\binom{2a}{a} 
%\binom{2\ell+1}{2k}
%\binom{\ell-k}{\ell-a}
%(-4)^{-a} 
%\left(
%\frac{t+1}{1-t}
%\right)
%^{2k}
\]
Additionally, from Lemma \ref{hodai1}, 
\[\begin{aligned}
G_{\ell, \ell}&=
4^{-\ell} 
(t+1)^{2\ell} (1-t)^{2\ell}
\sum_{k=0}^\ell
(-1)^{\ell-k}
\frac{2\ell+1}{2 \ell -2k+1}
\binom{\ell}{k} \binom{2\ell}{\ell}
\left(
\frac{t+1}{1-t}
\right)
^{2k}\\
&=
4^{-\ell} 
(t+1)^{2\ell} (1-t)^{2\ell}
\sum_{k=0}^\ell
(-1)^{k}
\frac{2\ell+1}{2k+1}
\binom{\ell}{k} \binom{2\ell}{\ell}
\left(
\frac{t+1}{1-t}
\right)
^{2\ell - 2k}.
\end{aligned}\]
Moreover, from Lemma~\ref{hodai2},
\[\begin{aligned}
\sum_{\ell \ge 0} G_{\ell, \ell} \lambda^\ell
&=
\sum_{\ell \ge 0} 4^{-\ell} 
(t+1)^{2\ell} (1-t)^{2\ell}\lambda^\ell
\sum_{k=0}^\ell
(-1)^{k}
\frac{2\ell+1}{2k+1}
\binom{\ell}{k} \binom{2\ell}{\ell}
\left(
\frac{t+1}{1-t}
\right)
^{2\ell - 2k} \\
&=
\sum_{k \ge 0}
(-1)^k
\left(
\frac{t+1}{1-t}
\right)
^{- 2k}
%
\sum_{\ell \ge 0} 
%
\frac{2\ell+1}{2k-1}
\binom{\ell}{k} \binom{2\ell}{\ell}
\left(
\frac{\lambda}{4}
(t+1)^4
\right)
^{\ell} \\
&=
\sum_{k \ge 0}
(-1)^k
\left(
\frac{t+1}{1-t}
\right)
^{- 2k}
%
\binom{2k}{k}
\frac{\left(
\frac{\lambda}{4}
(t+1)^4
\right)^k}{
\left(1 - (t+1)^4 \lambda \right)^{k+\frac{3}{2}}}.
\\
&=
\frac{1}{
\left(1 - (t+1)^4 \lambda \right)^{\frac{3}{2}}}
\sum_{k \ge 0}
\binom{2k}{k}
\left(
-
\left(
\frac{1-t}{t+1}
\right)^2
\frac{\frac{\lambda}{4}(t+1)^4}{1 - (t+1)^4 \lambda}
\right)^k\\
&=
\frac{1}{
\left(1 - (t+1)^4 \lambda \right)^{\frac{3}{2}}}
\sum_{k \ge 0}
\binom{2k}{k}
\left(
-\frac{\lambda (t+1)^2 (1-t)^2}{4(1 - (t+1)^4 \lambda)}
\right)^k\\
&=
\frac{1}{
\left(1 - (t+1)^4 \lambda \right)^{\frac{3}{2}}
\sqrt{1 +   \frac{ \lambda (t+1)^2 (1-t)^2}{1 - (t+1)^4 \lambda}}}\\
&=
\frac{1}{
(1- (t+1)^4 \lambda)\sqrt{
1-4 t (t+1)^2 \lambda
}},
\end{aligned}\]
as desired.
\end{proof}

\end{document}
